\section{Local passage laws from independent walk blocks}
\label{tree:sec:certificates}

The moment bounds give compactness of distances, but our later changes
of field require more: a conditional law for all paths confined to a
given open set. We construct that law in three steps. First retain
paths together with their ranges and ordered subpaths. Next reconstruct
confined paths from finitely many locally readable walk intervals.
Finally compare independent interval experiments with the full walk by
an explicit density and bridge calculation. Locality will mean dependence
of this conditional law on both $h$ and $\widehat h$ in the open set;
the passage object need not be a function of $h$. The final subsection,
Section~\ref{tree:sec:coherent-background}, retains one common continuum
background through changes of scale.

\subsection{Paths with ranges and ordered marks}

Write $\overline\R_+=[0,\infty]$,
metrized by $c\mapsto c/(1+c)$, with the usual convention at infinity.  For a
compact metric space $E$, write $\mathcal H(E)$ for its nonempty compact
subsets with the Hausdorff metric.  It is a compact metric space.  The empty
set can be adjoined as an isolated point when recording a collection that
might contain no paths.

Fix a compact time window $J$. At discrete resolution $n$, a path is a finite
sequence of primal edges, with zero-cost changes between contour
representatives of the same primal vertex also permitted.  Its cost is its
number of primal edges divided by the deterministic denominator $B_n$.
Here and throughout this section $B_n\to\infty$, and the interval estimates
of Lemma~\ref{tree:lem:moment_compactness} hold in these units.  We can use either
the whole sewn graph restricted to representatives in $J$, or one of the
finite instruction graphs defined below.  An edge in the latter is given
the same cost as the corresponding primal edge.
In an instruction graph, equality of vertices uses only its enabled
first-coordinate identifications. Two window copies are not identified
merely because their images coincide in a larger graph.

For every integer $m\geq1$ retain \emph{all} paths with $m+1$ ordered marks,
allowing repeated marks.  The data of a marked path are
\begin{equation}\label{tree:eq:marked_certificate_space}
 (t_0,\ldots,t_m;K_1,\ldots,K_m;c_1,\ldots,c_m)
 \in J^{m+1}\times\mathcal H(J)^m\times\overline\R_+^m.
\end{equation}
The $t_i$ are contour representatives of the marked vertices; $K_i$ contains
all representatives used on the portion between marks $i-1$ and $i$; and
$c_i$ is any upper bound for its cost. The closure of these data is the
set of \emph{path certificates}. If a vertex has several representatives, all choices are
retained, rather than choosing a preferred representative.  For a finite
instruction graph, replace $J$ by the finite disjoint union of its active
internal window copies, as defined below.
We also retain the ranges projected into a fixed compact coordinate
chart by the continuum curve in the contour coupling. The raw time
ranges and their spatial projections remain separate coordinates.

Deleting an interior mark unions the adjacent ranges and adds their
costs. Both operations are continuous in
Equation~\eqref{tree:eq:marked_certificate_space}, with addition in
$\overline\R_+$. Repeated marks may always be inserted. Hence the
deletion identities survive Hausdorff limits and retain traversal order,
which a compact range alone would not record.

There are three related collections of observations. The raw array
$\mathcal H_n$ records paths of the whole walk. A finite chart output
records paths using only the internal windows and cross-window links
declared by that chart; it is calculated from the same walk. Independent
reference experiments will instead sample the chart inputs afresh to
identify the conditional laws of these outputs. We retain their joint
\emph{laws} as auxiliary coordinates, while keeping $\mathcal H_n$
and its spatial projection throughout.

Extract jointly over bounded rational windows, finite rational cut lists,
all numbers of marks, and the finite instruction graphs defined below.
The Poisson cuts are independent of $\mathcal H_n$, the walk, and the
coupled background. The input marginals of the reference experiments
in Lemma~\ref{tree:lem:finite_kernel_transfer} converge by the invariance
principle and lattice approximation; their continuum laws stay fixed
during later extractions. Include beforehand the images of the instruction
lists under every quantum-unit transform used in the proof. These are
the countably many fixed dyadic transforms. Their band endpoints may
be nonrational, but remain a countable set of constants; rational bands
give the same exhausted collection by the reconstruction with open
margins below.
Retain internal endpoint distances and time diameter bounds as well.
Lemma~\ref{tree:lem:moment_compactness} makes these finite-valued variables
tight, while the hyperspaces are compact and the contours converge
locally uniformly. A diagonal extraction in their countable product
gives the \emph{certificate topology}.

The continuum background has its fixed joint law with the limiting
contour: fields, curve, reference metric, and residual embedding or
imaginary-field randomness are all retained. To test this coupling,
approximate the conditional expectation of a bounded background test
given the contour by bounded continuous contour functions in $L^1$ of
the fixed contour marginal. Use this approximation against bounded
\emph{raw} certificate tests, before spatial projection. The closure
lemma below and Equation~\eqref{tree:eq:bootstrap_background_closure}
give the details. This step needs no continuity of the map from
contours to fields or metrics.

\begin{proposition}[Certificate limit]\label{tree:prop:certificate_limit}
Suppose a sequence of bilateral sewing models satisfies the interval
modulus and moment bounds of Lemmas~\ref{tree:lem:moment_compactness}
and~\ref{tree:lem:moment_ui}, with $B_n\to\infty$.  There is a joint subsequential
certificate limit with the following properties, almost surely.
\begin{enumerate}
\item Every finite-cost certificate has approximating paths along the
selected subsequence.  Conversely, every sequence of paths with bounded
costs and with ranges in a fixed compact subset of an open coordinate
chart has a certificate subsequence.  The same assertions hold with any
fixed finite number of ordered subpaths marked.
\item If $K\Subset U$, endpoints converging to the same point of $K$ can
be joined at costs tending uniformly to zero by paths in $U$.  This holds
for discrete endpoints, limiting endpoints, and mixtures of the two in
the contour coupling.
\item Certificates concatenate finitely, with the sum of their costs
and with arbitrarily small additional cost in an arbitrary open
enlargement of their ranges.  They can be cut by marking approximating
paths and extracting further.
\item Let $F_U(x,y)$ be the infimum of the costs of certificates from
$x$ to $y$ with compact range in $U$.  On each connected open $U$ this
is a finite continuous pseudometric.  For $V\subset U$ one has
$F_U\leq F_V$ wherever both are defined.  All these assertions are
compatible with further certificate extractions.
\end{enumerate}
The assertion about continuity is local in $U\times U$; it does not assert
convergence of distances constrained to the boundary of a fixed closed set.
\end{proposition}

\begin{proof}
The main point is local joining; compactness alone does not give it.
On a fixed window let
$Z_n\to Z=(L,R)$ uniformly.  A first-coordinate identification in the
discrete sewing has equal endpoint heights and no smaller height between
them.  A second-coordinate match has this property up to one height step.
Consequently, every limit of a sequence of zero-cost identifications or
single edges is a Brownian identification.  Ordinary tree edges have
adjacent contour times and satisfy the same conclusion.  By compactness,
the largest projected diameter of a discrete link tends to zero.

By the Brownian fiber description in Section~\ref{tree:sec:background},
each nonsingleton class belongs to one coordinate and has two or three
times. At a triple $a<b<c$, the middle time is a strict local minimum,
with no further contact in $(a,b)\cup(b,c)$. We show that each such
identification can be approximated by a discrete link and short
internal paths.

Here is the required converse approximation.  Suppose
$f(a)=f(b)=H$ and $f>H$ on $(a,b)$, for $f=L$ or $R$.  Fix a point in
$(a,b)$.  At levels $H+\delta$, take the last crossing before that point
and the first crossing after it.  For sufficiently small $\delta>0$ the
crossings exist, the contour stays above that level between them, and
their times converge to $a,b$ as $\delta\downarrow0$.  To see the last
statement, for each $e>0$ the minimum on $[a+e,b-e]$ is strictly above
$H$.  First fix a level strictly below that minimum, and then use uniform
contour convergence.  The corresponding discrete excursion at a nearest
lattice level supplies a discrete identification, or a matched edge in
the second coordinate.  Its cost is at most $1/B_n$.

At $a<b<c$ with common height $H$, this construction above $H$ supplies
the pairs $(a,b)$ and $(b,c)$.  Below $H$ it supplies $(a,c)$: the two
outer times have values below $H$ arbitrarily close on their exterior
sides, since otherwise there would be another local minimum of height
$H$.  Taking the excursion above $H-\delta$ that contains $b$ gives the
claim.  Thus all three contacts, including the outer pair, are
approximable.  At each stage nearby representatives are joined by
short-time internal paths.  Their costs tend to zero by the interval
modulus, and their ranges stay in any prescribed neighborhood of the
contact by continuity of $\eta$.  The absence of mixed-coordinate
classes shows that this treats every surface identification.

\emph{Uniform joining.} The pointwise statement is uniform on compact sets.
Indeed, otherwise there would be $x_n,y_n$ with projected separation
tending to zero in a compact $K\Subset U$, for which the least joining
cost in $U$ stays above some $e>0$.  Properness of the bilateral curve
bounds all their time representatives in a common window.  After
subselection these representatives tend to $a,b$ with $\eta(a)=\eta(b)$.
The preceding chord construction, preceded and followed by short-time
paths, gives a joining path of cost tending to zero.  Its range tends
to that common point and is eventually in $U$, a contradiction.

\emph{Extraction and concatenation.} Hausdorff convergence of the closed certificate sets gives approximation
and subextraction in both directions.  A finite-cost limiting element
can be approximated by elements with finite costs converging to its
costs.  Elements in the discrete closure are themselves approximable by
actual paths.  If its range is compactly in $U$, both approximations
are eventually in $U$.  Marking a sequence of paths first and then using
the corresponding compact space proves the statement about subpaths.
The continuous deletion-of-marks maps ensure consistency with an already
retained unmarked certificate.

For concatenation, approximate a finite list of certificates.  Endpoints
that agree in the limit are connected by the uniform joining statement;
the connecting paths have total cost tending to zero.  The concatenated
paths have a further certificate limit with the required ranges and
cost.  Reversal is immediate.  For local finiteness, small time intervals
and the finite chord description give finite-cost local connections;
a finite chain of these connections along an ordinary continuous path
in $U$ connects any two prescribed points.  The same joining modulus,
the triangle inequality, and reversal show local continuity of $F_U$.
Monotonicity follows from its definition.  Every argument used compact
spaces, strict spatial margins, and the retained interval modulus, so
applies without change to a sequence of already-constructed limits.
\end{proof}

\subsection{Finite graphs determined by local cells}

We now describe finite graph tests that capture every compactly confined
path. Their inputs will be interval increments and relative frontier
heights, all readable from the fields in the confinement region.

Add a nested family of independent Poisson cuts in area time, with
intensities $2^q$, $q\geq0$, by independent superposition.  Their images
form the corresponding nested Poisson point processes of intensities
$2^q\mu_h$.  We call the image of a consecutive cut interval a cell.
All cells used below have compact containment in the stated open set.
Their orientation, area parametrization, and incoming and outgoing
boundary arcs are part of the local data of
Proposition~\ref{tree:prop:background_locality}; see also
\cite[Lemma~2.4]{GwynneMillerSheffieldHarmonic} and \cite[Lemma~3.7]{ContrerasHipGwynneCurvature}.
The abstract quantum surface of a run of consecutive cells is determined
by its two contour increments
\cite[Lemma~4.9]{GHSSurvey}.  The embedded run is read from the local cells.

\paragraph{Reading the interval increments.}
For an oriented subsegment $C_{[a,t]}=\eta([a,t])$ in the retained
cells and $X=L$ or $R$, put $m_X=\min_{[a,t]}X$.  Its tagged incoming
and outgoing $X$-frontiers have respective quantum lengths
$X(a)-m_X$ and $X(t)-m_X$.  These are lengths on the abstract interval
surface, with frontier multiplicity; overlapping geometric arcs are
not replaced by their union.  Consequently
\begin{equation}\label{tree:eq:local_frontier_increment}
 X(t)-X(a)=\nu_h(\partial_X^{\mathrm{out}}C_{[a,t]})
                 -\nu_h(\partial_X^{\mathrm{in}}C_{[a,t]}).
\end{equation}
The increment identity is the boundary-length definition in
\cite[Equation~(4.3), p.~35, and Figure~8, p.~36]{GHSSurvey};
the four boundary segments are described in
\cite[discussion following the proof of Corollary~9.3, p.~152]{DMS}.
The locally determined oriented curve and its stopped left and right
flow lines distinguish the two sides and the incoming and outgoing
frontiers of this segment. No order of other disconnected segments is
needed. Relative time is local quantum area. To read the lengths, use
the fact that the quantum
length of an $\SLE_2$ frontier is the fixed-normalization
$(\gamma/2)$-Gaussian multiplicative chaos of its $5/4$-dimensional
Minkowski-content measure; see
\cite[Definition~2.11 and Proposition~3.3]{BenoistLength} and the
local absolute-continuity formulation in
\cite[paragraph following Theorem~4.1, p.~33]{GHSSurvey}.
On a compact frontier subarc in $U$, Minkowski content is determined
by the arc and sufficiently small field mollifiers use only $h|_U$.
The chaos limit, followed by exhaustion, is therefore locally
measurable. Construct these measures simultaneously on the countable
imaginary-geometry branches generating the curve \emph{before} selecting
subarcs by Poisson area cuts or other real-field data. Then restrict
the measures to the selected subarcs, keeping each tagged frontier
occurrence separately. This preserves multiplicity when a geometric
arc appears on more than one frontier, and requires no independence
between a selected subarc and the real field.
Equation~\eqref{tree:eq:local_frontier_increment}
therefore recovers the increments at every rational relative area time;
continuity recovers the whole two increment paths.  Absolute contour
height and absolute global time cancel in the difference.

\paragraph{A chart on whole runs.}
We first describe the graph when each retained run is used as one whole
block. A \emph{reference block}
is an oriented interval
$I_b=[0,\ell_b]$ together with its two increment functions
$z_b^L,z_b^R$, starting at zero, and a finite set of internal cuts.
There are finitely many blocks in a package. For each coordinate
$a\in\{L,R\}$ choose a forest on the blocks and give its edges offsets
$o_{bd}^a=s_d^a-s_b^a$.  A choice of zero at one block in each forest
component determines all $s_b^a$ in that component.  No expression below
uses heights from different components.  A component-wide translation
has no effect.

Give each block its own copy of the contour vertices. Inside that copy,
enable the ordinary sewing witnessed in the block. Between different
copies, enable only the band instructions defined next. Thus the forest
offsets are the only independently supplied relative heights; the
internal strings and declared bands determine every available link.

Put $m_b^a=\min z_b^a$ and $e_b^a=z_b^a(\ell_b)$.  The incoming arc of
block $b$ in coordinate $a$ is parametrized by the levels of the
left-record set, and its outgoing arc by those of the right-record set:
\begin{align}\label{tree:eq:chart_arcs}
 \mathcal I_b^a&=\{t:z_b^a(t)=\min_{0\leq u\leq t}z_b^a(u)\},&
 \mathcal O_b^a&=\{t:z_b^a(t)=\min_{t\leq u\leq\ell_b}z_b^a(u)\}.
\end{align}
For this \emph{continuum boundary parametrization}, repeated scalar
representatives at a level refer to one boundary point.  This statement
does not identify primal vertices in the discrete second coordinate.
The
height ranges of these arcs are respectively $[m_b^a,0]$ and
$[m_b^a,e_b^a]$.

A band instruction consists of an ordered pair $(b,d)$, a coordinate
$a$, two rational numbers $\alpha<\beta$, and a finite list $W$ of
blocks in the same offset component.  It permits outgoing-to-incoming
links at all levels $H$ such that
\begin{equation}\label{tree:eq:band_instruction}
 \alpha\leq H-s_b^a\leq\beta,
 \qquad H\leq s_w^a+m_w^a\quad(w\in W).
\end{equation}
The outgoing and incoming arc conditions in
Equation~\eqref{tree:eq:chart_arcs} are always imposed.  The empty list $W$
is allowed. Ordinary sewing is enabled only inside each block of this
test. Cross-block first-coordinate links
identify vertices and have cost zero.  The second-coordinate rule is
instead a rule for \emph{oriented step incidences}.  In unscaled integer
heights an outgoing incidence is an up-step with terminal index $u$,
lower height $h=R_{u-1}$, upper height $h+1=R_u$, and
$\min_{[u,\,\text{end of block}]}R\geq h+1$.  An incoming incidence
is a down-step with terminal index $v$, with $R_{v-1}=h+1$, $R_v=h$,
and $\min_{[\text{start of block},\,v-1]}R\geq h+1$.
Test the entire strip $[h,h+1]$ against the band and minimum comparisons.
A permitted pair adds an edge from the primal vertex $[u-1]_L$ to
$[v-1]_L$, of cost $1/B_n$.  There is no zero-cost second-coordinate
quotient, even for repeated visits to the same record level.  This
retains every edge cost on paths with arbitrarily many matches.
Scaled heights give the identical rule with strip width $n^{-1/2}$.
If a comparison is attained at a selected block minimum, use the
exact discrete minimum, rather than a rounded continuum value.

\paragraph{Retaining old tests when runs are merged.}
The whole-block graphs suffice for capturing paths. Joint tests and
refinements require one additional distinction. Suppose two old runs
$I_1$ and $I_2$, separated by an omitted gap, are included in a larger
retained interval. The larger string determines their relative starts
from its increments. But the old test still permits internal edges
only when witnessed in $I_1$ or $I_2$, together with its declared bands;
a new graph on the larger interval also permits edges witnessed in the
gap. Restricting the ranges of paths in that larger graph does not
recover the old graph. We therefore keep both graphs as separate
outputs and record the old intervals as windows of the larger string.

In the general package, call each supplied string a \emph{master block}.
A test selects a finite list of \emph{active internal windows}, each
an interval $[u,w]$ between recorded cuts of a master block $M$, with
its own copy of the contour vertices. Its starting height is
$s_M^a+z_M^a(u)$ and its increment path is
$z_M^a(u+\cdot)-z_M^a(u)$. These starts are derived data, not new
independent forest offsets. Enable ordinary sewing only as witnessed
inside that window, and use the preceding band rules for cross-window
links, with the windows' derived heights and increments. The simple
case above has one whole active window per master block.

To preserve every old test under refinement, an arc or
minimum may be taken on any subinterval whose endpoints are recorded
cuts of its block.  A band may use any recorded cut as its height
anchor.  If the anchor is time $u$ in block $b_0$, replace $s_b^a$ in
the band inequality by
$q_A^a=s_{b_0}^a+z_{b_0}^a(u)$; thus its two endpoints are
$q_A^a+\alpha$ and $q_A^a+\beta$.  A minimum on a recorded window
$[u,w]$ in block $b_1$ means exactly
$s_{b_1}^a+\min_{[u,w]}z_{b_1}^a$.  This adds only finite selectors
from the internal cut lists to the instruction template.  The anchor
heights and restricted minima are evaluated from the same local
increments; they are not additional independently sampled offsets.
Equation~\eqref{tree:eq:band_instruction} is the case of a start anchor and
whole-block windows.  This closure is needed because, after merging
runs, the start of an old run is an internal cut of the merged run.
It also lets a band retain its exact height endpoints upon refinement,
instead of rounding an increment-dependent endpoint anew.

The formal graph uses only these instructions, its active-window list,
the master block strings, and
the forest offsets.  In particular it does not query any omitted portion
of the chronological walk.  We retain its complete marked certificate
sets, including confinement data in the constituent blocks.  There are
only countably many instruction lists: durations, increments, and offsets
are input variables, not names in the list.  Global time locations and
the chronological order of blocks that have no declared seam are not
inputs to an instruction graph.

For a fixed instruction list, place each active window on its own tagged
copy of $[0,1]$, using relative time divided by that window's duration.
Retain the durations and cuts among the inputs. The marked output
hyperspaces are then fixed compact spaces even when the input durations
vary. Actual time-labelled ranges are recovered from these inputs;
this is only a coordinate convention for the finite chart outputs.

\paragraph{Admissible bands.}
We next specify which instructions can be recognized from embedded cells.
A directed shared seam is a common subarc of positive boundary length
of an outgoing arc and an incoming arc in the same coordinate.  On this
arc the two height parameters differ by a constant.  That constant is
the offset; it is read at any two interior points, so choices of an
endpoint are unnecessary.  In testing admissibility, first clip the
rational band to the two natural record-arc ranges.  A band with no
cap is admissible if this clipped interval lies in the shared seam
and has positive margin from its upper end.  At its lower end it
must either have positive margin or end at a natural record-arc
minimum of a selected block, with a positive shared interval immediately
above that minimum.  In the latter case the rational lower endpoint
is chosen strictly below the minimum: the exact discrete record-arc
condition performs the clipping, at every resolution.  A band may
extend across the upper end $H$ of a seam
provided that its cap is the exact minimum of a selected block $w$
and the three incident directed seams are visible: the outer pair
overlaps below $H$, and each outer block overlaps $w$ above $H$.
All other effective band endpoints must have positive margin.  We allow several
such caps; a cap that is strictly larger than the whole band is harmless.
Bands can also be restricted to a compact subarc of their constituent
cells when confinement is tested.  Call these \emph{buffered admissible
charts}.  All parts of this definition are measurable from the embedded
oriented cells and their boundary length parametrizations.  Choosing a
rational interior level, an integer lower bound on the margins, and a
spanning forest can be done from a countable list.

\begin{lemma}[Finite-band reconstruction]\label{tree:lem:chart_reconstruction}
Almost surely the following holds simultaneously for all compact
$K\Subset U$ in a countable basis of bounded coordinate patches and,
by exhaustion, for all compact $K\Subset U$.
\begin{enumerate}
\item Sufficiently fine selected cells in $U$ form finite buffered
admissible charts that capture every sequence of discrete paths whose
projected ranges approach $K$.  A single finite collection of whole
bands suffices for all the links of all these paths, regardless of
their number of seam crossings.
\item Conversely, a bounded-cost chart path whose limiting range has
compact containment in $U$ is realized by an actual sewn-graph path
in any open enlargement of that range.  Its cost is unchanged, apart
from vanishing endpoint joining costs.
\item Taking common refinements, unions of finitely many charts, and
closure with arbitrary open margins recovers exactly the certificates
of Proposition~\ref{tree:prop:certificate_limit} that are compactly in $U$.
The statement includes all finite marked-subpath data.  All forest
offsets and chart selections required for this reconstruction are
measurable from the fields and Poisson marks in $U$.
\end{enumerate}
The same conclusions hold inside the interior of a fixed interval cell,
using only links witnessed within that interval.
\end{lemma}

\begin{proof}
The proof has three steps: identify which gap minima obstruct a link,
read every necessary cutoff from local seams, and retain finitely many
whole bands. The last step retains every crossing edge without an
approximation error for each crossing.

\emph{The scalar test and omitted gaps.} Let
$I_b=[a_b,b_b]$ and $I_d=[a_d,b_d]$ be disjoint chronologically ordered
runs, with $b_b<a_d$.  A chord from a time in $I_b$ to a time in $I_d$
at height $H$ exists precisely when the first time is on the outgoing
record arc, the second is on the incoming record arc, their heights
are $H$, and
\begin{equation}\label{tree:eq:gap_minimum_condition}
 H\leq \min_{b_b\leq t\leq a_d}f(t).
\end{equation}
Indeed these conditions say exactly that $f\geq H$ between the two
times.  Splitting the intervening interval into selected runs and omitted
gaps splits its minimum into the minima of those finitely many pieces.
This is also an exact discrete statement, with the full-step convention
for an edge.  Thus a whole overlap, in height coordinates, is an interval
with upper endpoint given by one of these finitely many minima, further
intersected with the two record-arc ranges.  This calculation proves
both the sufficiency and the necessity of the band language.

Fix a positive $r$ with the closed $3r$-neighborhood of $K$ compactly
in $U$.  Properness bounds its time preimage in a finite window.  Choose
a cut refinement for which every cell intersecting that neighborhood
has diameter less than $r$ and belongs to a slightly larger bounded
time window.  This is possible by uniform continuity of $\eta$ and
the fact that the largest Poisson gap in a fixed window tends to zero.
Select all cells meeting the closed $2r$-neighborhood of $K$.
They are finitely many and compactly in $U$.  Merge consecutive selected
cells into maximal runs.  Every visit to a point in the $r$-neighborhood
of $K$ lies in a selected cell.  The endpoints of a maximal run stay
away from $K$: if such an endpoint were close to $K$, the adjacent cell
would also have been selected.  Independent cut times are not nontrivial
Brownian chord times, by conditioning on the cuts and using the
Brownian facts in the proof of Proposition~\ref{tree:prop:certificate_limit}.

Consider all possible limits of links whose projections approach $K$.
Their endpoint time pairs form a compact set.  If such a pair determines
a chord at level $H$, then \emph{every} intervening time attaining $H$
represents the same surface point.  An omitted gap cannot have minimum
$H$: a cell containing a time attaining that minimum would then have
been selected.  Hence every relevant omitted-gap minimum is strictly
larger than $H$.  There is a uniform positive difference after fixing
the finite runs and restricting to the compact set of relevant link
pairs.  Otherwise a sequence of pairs and one of the finitely many
omitted gaps would converge to a chord with its minimum attained in
that omitted gap.  The preceding observation rules this out.  This
argument also covers minima at gap endpoints, since their two adjacent
cells would have been selected.

It follows that omitted gaps can be removed from
Equation~\eqref{tree:eq:gap_minimum_condition} for all the relevant levels,
provided bands are chosen with strictly positive slack below their
minima.  A cutoff that remains is the exact minimum of a selected
intervening run.  Retaining its comparison in
Equation~\eqref{tree:eq:band_instruction} implements that cutoff exactly
at every discrete resolution.  This exact comparison matters: replacing
it by an approximate real cutoff could add a spurious edge at a branch.

\emph{Reading offsets and exact cutoffs.}
Away from an endpoint of an overlap, a compact band in a
directed shared seam reads the difference $s_d^a-s_b^a$ from the two
boundary length coordinates.  Suppose a cutoff $H$ is attained in a
selected run $w$.  There are three times $t_1<t_2<t_3$ representing the
chord point, with $t_2\in I_w$.  The scalar class description shows
that $t_2$ is the sole intervening contact and a strict local minimum.
The cut times lie away from all three visits.  Choose small neighborhoods
of the visits inside their respective runs.  On compact subintervals
of $(t_1,t_2)$ and $(t_2,t_3)$ the contour is strictly above $H$.
The chord approximation in Proposition~\ref{tree:prop:certificate_limit}
then gives positive intervals of levels above $H$ shared by the first
and middle runs and by the middle and last runs.  Immediately below
$H$, the two outer runs share an interval of levels.  The directions
of these three overlaps distinguish the middle run.  They recover its
offset to the outer runs, and its recorded internal minimum supplies
the cap.  This proves local recognizability even though the branch
point itself has zero boundary length.
See Figure~\ref{tree:fig:triple-contact}.

\begin{figure}[tbp]
\centering
\begin{tikzpicture}[x=1.65cm,y=1.3cm,>=Stealth,font=\small]
  \fill[blue!7] (0.45,-0.35) rectangle (5.55,0);
  \fill[orange!12] (1,0) rectangle (3,0.35);
  \fill[orange!12] (3,0) rectangle (5,0.35);
  \draw[->,gray] (0,-0.98)--(0,1.9) node[above]{height};
  \draw[->,gray] (0,-0.95)--(5.95,-0.95) node[right]{time};
  \draw[gray,dotted] (0,0)--(5.6,0);
  \node[left] at (0,0) {$H$};
  \node[left] at (0,-0.3) {$H-\varepsilon$};
  \node[left] at (0,0.3) {$H+\varepsilon$};
  \draw[very thick] (0.45,-0.55)--(1,0)
    .. controls (1.35,1.2) and (1.15,1.55) .. (1.65,1.35)
    .. controls (2.05,1.1) and (1.85,0.6) .. (2.2,1.0)
    .. controls (2.6,1.2) and (2.55,0.25) .. (3,0)
    .. controls (3.3,0.15) and (3.2,1.0) .. (3.65,0.8)
    .. controls (4.0,0.55) and (4.0,1.65) .. (4.35,1.5)
    .. controls (4.7,1.2) and (4.72,0.55) .. (5,0)
    --(5.55,-0.55);
  \draw[blue!70!black,thick,<->] (0.7,-0.3)--(5.3,-0.3);
  \node[blue!70!black,fill=blue!7,inner sep=1pt] at (3,-0.3)
    {outer seam};
  \draw[orange!75!black,thick,<->] (1.1,0.3)--(2.65,0.3);
  \draw[orange!75!black,thick,<->] (3.25,0.3)--(4.85,0.3);
  \node[orange!75!black] at (1.9,1.85) {adjacent seam};
  \node[orange!75!black] at (4.15,1.85) {adjacent seam};
  \foreach \x/\i in {1/1,3/2,5/3}{
    \fill (\x,0) circle (1.8pt);
    \node[fill=white,inner sep=1pt] at (\x,-0.65) {$t_\i$};
  }
\end{tikzpicture}
\caption{A triple contact in one contour coordinate. The visits
$t_1<t_2<t_3$ at height $H$ represent the same surface point.
Overlaps below $H$ join the outer visits; overlaps above $H$ join
adjacent visits. Their directions determine the relative offsets and
identify the middle run, whose exact minimum supplies the cutoff $H$.
Omitted gaps stay separated from the retained bands by a positive
height margin. This diagram shows the incidences schematically.}
\label{tree:fig:triple-contact}
\end{figure}


The same configuration explains a lower endpoint that must not be
discarded.  A seam between an outer run and the middle run begins
at $H=m_w^a+s_w^a$ and has positive length above $H$.  Choose its
rational band to extend below $H$.  The exact incoming or outgoing
record condition in the middle run clips it at the discrete minimum
$s_{w,n}^a+m_{w,n}^a$.  Hence links at that minimum, as well as links
at lattice heights converging to it from above, are retained.  No
new comparison with an omitted gap is introduced.  Positivity of
the shared interval above $H$ already provides slack from every
upper cutoff.  More generally the lower endpoint of a shared seam
is the larger of its two natural record-arc minima, by
Equation~\eqref{tree:eq:gap_minimum_condition}; these exact natural
conditions handle it in precisely the same way.  Thus neither
lower endpoints nor upper branch cutoffs require an additional
joining cost for each crossing.

For each coordinate, join run vertices whenever such a positive
overlap is used and retain a spanning forest of this graph.  Differences
along the forest recover all offsets needed by its bands and caps.
Cycles require no additional random parameters: their sums are zero
because they are differences of the same height coordinates.  An
isolated component requires no absolute height and no comparison with
another component.  In particular there is no need to infer the
chronological ordering of two remote runs without a shared seam.
The direction of a declared seam specifies the order of its two ends;
the three directed overlaps specify an intermediate cap.  The omitted
chronological ordering is used only in the proof that the local
instructions are valid, and is not an input to them.

\emph{A finite family of whole bands.}
There are finitely many run pairs and gap pieces. For each relevant
pair, its height set is compact.  Cover it by bands with rational
endpoints lying in positive seam overlaps, allowing rational lower
endpoints below a natural minimum and clipping by the exact record-arc
condition as just explained.  Use the exact selected-minimum cap when
an upper overlap endpoint belongs to the height set.
The preceding positive-overlap and omitted-gap arguments supply a
neighborhood of every height, so a finite subcover exists.  Include
the finitely many additional overlaps used to recover the offset
forest.  This produces one finite instruction list.  Each instruction
retains \emph{all} discrete links in its band; it does not prescribe
one crossing of that band.  Consequently the same list captures paths
with an arbitrary, even diverging, number of crossings.

Uniform contour convergence now proves capture.  If a discrete link
near $K$ failed to be captured infinitely often, a limiting endpoint
pair would belong to the compact set just covered.  In its selected
band, outgoing and incoming record conditions are the original
discrete conditions, all selected minima are tested exactly, and
omitted gaps have a fixed positive slack.  The link therefore satisfies
the instruction for all sufficiently large resolutions, a contradiction.
Internal tree steps and internal matches are already included.

\emph{Realization and refinement.}
For realization, consider an admissible chart and a compactly confined
limiting path.  On a band strictly interior to a genuine shared seam,
Equation~\eqref{tree:eq:gap_minimum_condition} has positive margin from
every omitted obstruction.  At a cap the local three-overlap test
identifies the selected intermediate minimum; its exact discrete
comparison supplies the only possibly non-strict condition.  The
remaining omitted obstructions have positive slack on the confined
compact set, by the same minimum argument.  Thus each permitted link
of the discrete approximants is an actual link once the resolution
is sufficiently large.  This conclusion is uniform over all their
links, because there are finitely many bands and the slack is uniform.
If endpoint representatives differ, the joining conclusion of
Proposition~\ref{tree:prop:certificate_limit} corrects them in an arbitrary
open margin at vanishing cost.

A refinement retains all old cut anchors and active windows and adds
cuts.  An old active window remains the same whole cut-to-cut interval;
its internally enabled edge set is not changed by the extra cuts.
When master strings are merged, the old windows remain selectors in
the merged string, their starting heights are calculated from its
increments, and their band anchors retain their exact values.  Each
old test graph is therefore \emph{calculated directly} from the merged
input using its old edge-witness windows and band list.  It is retained
as its own output component, alongside any new graph enabled on the
larger union.  It is not recovered by restricting ranges of paths in
the larger graph, which could use newly enabled gap edges.  Additional
arc or minimum windows may be subdivided using the finite selector
language without changing the old active internal windows.  Thus
every old path and every ordered marking remains exactly an output
of its old test.  Countably many
such operations were included in the extraction.  Applying capture
and realization with margins decreasing to zero proves the third
assertion.  Measurability follows by using the countable instruction
lists and rational positive margins in the local tests above.

Finally, if $K$ is compactly in the interior of an interval cell
$\eta([a,b])$, all its curve-time preimages lie in $(a,b)$.  Indeed
interval interiors are disjoint from the images of disjoint time
intervals; this is also a consequence of the area parametrization
and the disk-cell property of Theorem~\ref{tree:thm:background_peanosphere}.
Compactness makes the containment uniform.  All cells selected near
$K$ are then internal to $[a,b]$, so the entire construction uses only
links witnessed in that interval.
\end{proof}

\subsection{Conditional transfer from independent intervals}

The finite graph tests use local data, but their limiting conditional
laws still need to be identified. We compare independent input blocks
with the same blocks in the full walk. Poissonizing the jump clock
removes lattice parity; the following lemma controls the interval
endpoints when we return to the original clock.

\begin{lemma}[Stability at nonidentified cuts]\label{tree:lem:cut_stability}
Suppose two pairs of discrete time endpoints converge to the same
$a<b$ in a common bounded contour window, and the Brownian fibers
of $a$ and $b$ are singletons.  Under the uniform internal time
modulus of Lemma~\ref{tree:lem:moment_compactness}, their internal
endpoint distances, in the common cost units, differ by
$o_{\Prob}(1)$.  If their $k$th powers are uniformly integrable,
their $k$th moments have the same limit.  This applies simultaneously
to every fixed finite set of deterministic or independent Poisson cuts.
\end{lemma}

\begin{proof}
It suffices to compare the smaller and larger intervals formed from
the two pairs of endpoints.  The added portions are shrinking time
fringes near $a$ and $b$.  Every edge available in the larger internal
graph but not the smaller one has a witnessing step or match endpoint
in one of these fringes.  Likewise an old vertex represented in an
added fringe is linked by a first-coordinate identification to that
representative.  Uniform contour convergence and singleton fibers
imply that all these incidences and their other representatives lie
in time neighborhoods of the same cut whose radii tend to zero in
probability.  Indeed, otherwise a subsequence in the containing
compact window would give a link with one limiting endpoint $a$
or $b$ and its other endpoint a different time, contrary to the
singleton assumption.  The two neighborhoods are disjoint eventually.

Take a path in the larger graph between its tips.  Every newly
available edge is confined to one of these two neighborhoods.  Keep
the part after the last visit to the left neighborhood before the
first visit to the right neighborhood.  Its middle uses only edges
of the smaller graph.  Its first and last old vertices have smaller
interval representatives in the respective shrinking cut neighborhoods.
Connect them internally to the smaller interval's tips using the
time modulus.  This adds $o_{\Prob}(1)$ once at each end.  It does
not pay separately for the possibly many discarded fringe visits.
Taking infima bounds the smaller interval's endpoint cost by the
larger one's plus this error.  Graph inclusion and short endpoint
pieces in the larger interval give the reverse inequality.  The
moment assertion follows from uniform integrability.  Fixed and
independent Poisson times have singleton fibers almost surely by
the Brownian facts used in Proposition~\ref{tree:prop:certificate_limit},
so conditioning on their countable collection proves the last claim.
\end{proof}

Put the bilateral axial walk on an independent Poisson jump clock of
rate $n$ in macroscopic time.
The coordinate walks are then independent continuous-time simple walks,
each of variance rate $\sigma^2=1/2$ before diffusive rescaling.  We write
$\epsilon_n=n^{-1/2}$ and use rescaled heights in
$\epsilon_n\Z$.  This removes parity constraints.  The clock satisfies
the uniform law of large numbers on each bounded time interval.
Replacing every jump time by its jump index is an \emph{exact}
relabeling of the full sewn graph.  A formal test cut at clock time
$T$ becomes a cut at index time $J_n(T)/n$; its active edge-witness
windows, oriented step incidences, height anchors, and costs are
identical.  The law of large numbers says only that these relabeled
cuts have the same limiting marks.  We do not infer continuity of
arbitrary formal test outputs under independently moved cuts.
For internal endpoint norms with fixed macroscopic endpoints,
Lemma~\ref{tree:lem:cut_stability} supplies the additional comparison.
The moment argument of Lemma~\ref{tree:lem:moment_ui} also applies in
Poisson time: disjoint clock intervals have independent strings,
their step counts concentrate, and the probability of a count
outside a fixed multiple of its mean decays exponentially.  The
trivial cost bound by the count controls this last event.  Thus
the same mixed-monomial argument gives uniform integrability, and
the cut-stability lemma transfers the internal $k$th moments as
well.  This justifies returning to fixed walk lengths.

\begin{lemma}[Finite reference experiment and conditional transfer]
\label{tree:lem:finite_kernel_transfer}
Fix finitely many chart packages.  A package may contain several
overlapping instruction tests; tests in different packages use disjoint
blocks and disjoint offset forests.  Give the block durations and
relative internal cuts a smooth positive density on their ordered
parameter domain.  Sample independent Poissonized increment strings
on the blocks. Independently, sample the forest offsets with smooth
positive probability densities, assigning each point of
$\epsilon_n\Z$ the density's mass in its centered half-open lattice
cell. Retain the complete
marked certificate data of the package's canonical instruction graphs.

There is a common subsequence on which these reference experiments,
for all finite packages in the countable instruction language, have
conditional limiting kernels
\begin{equation}\label{tree:eq:reference_kernel}
 \nu_r(dx_r)\,K_r(x_r,dy_r).
\end{equation}
Here $x_r$ consists exactly of the relative block durations, internal
cuts, increment paths, and forest offsets, and $y_r$ is the package's
additional certificate data.  The following statements hold.
\begin{enumerate}
\item Insert the selected blocks, in any chronological order consistent
with the desired actual word, with positive-duration omitted gaps.
Then in the limiting actual-walk experiment, conditional on the
\emph{entire} limiting contour and all cut marks, the law of the
additional package data is $\bigotimes_r K_r(x_r,\cdot)$.
\item The kernel is independent of absolute time locations, integer
cut labels, and the chronological ordering of pairs of blocks not
used by the instruction list.  It is unchanged by replacing the
smooth positive reference densities.  Relabeling blocks or replacing
an offset forest by another basis for the same differences transforms
the kernel by the corresponding deterministic relabeling of its inputs
and outputs.
\item These assertions hold jointly for finite collections of tests,
with their common refinements, and remain true after measurable
selection from the jointly retained countable family of instruction
templates and finite cut-mark lists. The selection may depend on the
limiting contour and marks. Selection is
performed after conditioning; no claim of this kind is made for
selection based on the additional passage data.
\end{enumerate}
\end{lemma}

\begin{proof}
We compute the density of the actual block offsets, insert the omitted
intervals as bridges, and then condition on the whole contour. We finish
by checking relabeling and refinement, which are needed for local chart
selection.

\emph{Independent reference experiments.} Under the
reference law, distinct packages are independent. Their input
marginals converge by the invariance principle and lattice
approximation of the offsets: sample a real offset from its density
and replace it by the center of its lattice cell, moving it by at most
$\epsilon_n/2$. Their output coordinates belong to
compact hyperspaces, so after extraction the joint laws converge to
Equation~\eqref{tree:eq:reference_kernel}.  The joint limit of independently
sampled packages is their product; this statement concerns joint weak
convergence, and imposes no regularity on the conditional kernels $K_r$.

\emph{The exact change of density.}
First treat one coordinate and fix the continuous cut parameters.
List all selected blocks in their actual chronological order
$1,\ldots,m$.  Let $s_i$ be their starting heights, $e_i$ their net
increments, and $g_i>0$ the gap duration from the end of block $i$ to
the beginning of block $i+1$.  Set $s_1=0$, since all graphs are
invariant under a common height translation.  Complete the union of
the package forests to a tree $T$ on the $m$ blocks.  Add independent
offset variables on its extra edges, with positive smooth densities.
For an edge $e$ with smooth positive probability density $\rho_e$,
its discretized mass is
\begin{equation}\label{tree:eq:offset_mass}
 r_{n,e}(a)=\int_{a-\epsilon_n/2}^{a+\epsilon_n/2}\rho_e(u)\,du,
 \quad a\in\epsilon_n\Z.
\end{equation}
The half-open cells partition $\R$, so these positive masses sum
exactly to one. Uniform continuity on compact sets gives
$\epsilon_n^{-1}r_{n,e}(a)\to\rho_e(a)$ uniformly for $a$ in a compact
set. Gaussian densities are one possible fixed choice; the construction
applies to every smooth positive probability density.

The map from the $m-1$ tree differences to
$(s_2-s_1,\ldots,s_m-s_{m-1})$ is unimodular.  To prove this, recover
each $s_i$ by summing signed edge differences along the unique path
in $T$ from $1$ to $i$.  This is an integer linear map.  Its inverse
is obtained by reconstructing the $s_i$ from the consecutive differences
and subtracting the appropriate two values for every tree edge.
Both maps have integer entries, so their determinants are integers
whose product is one.  Thus their determinants are $+1$ or $-1$,
and no extra lattice multiplicity occurs.

Let $p_t(k)$ be the transition mass of the unscaled coordinate walk
in time $t$.  Conditional on the selected block paths, the required
gap displacement is
$\epsilon_n^{-1}(s_{i+1}-s_i-e_i)$.  The exact density of the actual
relative starts with respect to the reference offset law is
\begin{equation}\label{tree:eq:forest_density}
 W_n^{a}=
 \frac{\displaystyle\prod_{i=1}^{m-1}
   p_{ng_i}\bigl(\epsilon_n^{-1}(s_{i+1}-s_i-e_i)\bigr)}
 {\displaystyle\prod_{e\in T}r_{n,e}(a_e)}.
\end{equation}
Indeed, summing its numerator over all consecutive gap displacements
gives one, and unimodularity identifies this sum with the sum over
tree offsets.  Multiply Equation~\eqref{tree:eq:forest_density} for the
two coordinates and multiply by the ratio of cut-parameter densities.
This is the full density $W_n$.  If an absolute starting value is
needed for reconstructing the whole contour, include an anchored
additional block, or record that value separately; no chart graph
uses it.  External past and future portions are added conditionally
on the extreme endpoints in the same way as the gaps.

\emph{Uniform density and bridge limits.}
The Fourier formula for a continuous-time coordinate walk is
\begin{equation}\label{tree:eq:poisson_walk_fourier}
 p_{nt}(k)=\frac1{2\pi}\int_{-\pi}^{\pi}
 \exp\{-n\sigma^2t(1-\cos\theta)-ik\theta\}\,d\theta.
\end{equation}
After $\theta=\epsilon_n u$, the integrand converges to
$\exp\{-\sigma^2tu^2/2-ixu\}$ when
$\epsilon_n k\to x$.  The bound
$1-\cos\theta\geq2\theta^2/\pi^2$ for $|\theta|\leq\pi$
gives an integrable Gaussian majorant uniformly when
$t\in[t_-,t_+]\subset(0,\infty)$.  Truncating the $u$ integral
first also gives uniformity in $x$ on compact sets.  Therefore
\begin{equation}\label{tree:eq:uniform_gap_lclt}
 \sup_{\substack{t\in[t_-,t_+]\\ |\epsilon_n k|\leq M}}
 \left|\epsilon_n^{-1}p_{nt}(k)-
 \frac{e^{-(\epsilon_n k)^2/(2\sigma^2t)}}
 {\sqrt{2\pi\sigma^2t}}\right|\longrightarrow0.
\end{equation}
There are $m-1$ numerator and denominator lattice factors in
Equation~\eqref{tree:eq:forest_density}; they cancel.  On compact parameter
sets with positive gap durations, $W_n$ thus converges uniformly to
a finite positive continuous function $W$.

Conditional on the selected block paths and their starts, the omitted
gaps are independent bridges and are independent of every $y_r$.
Their rescaled laws converge uniformly, on the same compact parameter
sets, to the corresponding Brownian bridges.  Here is a proof of the
uniform assertion.  Their finite-dimensional densities are products
of the transition masses divided by the whole-gap mass; applying
Equation~\eqref{tree:eq:uniform_gap_lclt} proves convergence at times
bounded away from each other and the endpoints.  On the first half
of a gap of duration $t$, the bridge density with respect to the
unconditioned walk started at its first endpoint is
\[
 \frac{p_{n(t-u)}(k-X_{nu})}{p_{nt}(k)},\qquad u\leq t/2.
\]
The denominator is at least $c\epsilon_n$ for compact endpoint
displacements and $t\in[t_-,t_+]$.  The Fourier bound gives a numerator
at most $C\epsilon_n$ uniformly in its integer argument.  Thus the
density is uniformly bounded.  The free-walk invariance principle
gives tightness on this half.  Time reversal proves it on the second
half.  Together they prove bridge tightness; the finite-dimensional
limits identify every subsequential limit, proving uniform convergence.
The same argument applies to a finite vector of gaps and to external
contour pieces on any bounded window.

\emph{Conditioning on the complete contour.}
Let $a$ denote the added completion offsets and let $z$ denote all
omitted bridge paths.  Test the joint law against products of bounded
continuous functions of the inputs, outputs, and bridges.  Integrate
the bridge variables first.  On a compact parameter set, uniform
bridge convergence replaces the bridge conditional expectation by
its Brownian counterpart, and uniform convergence of $W_n$ replaces
the density by $W$.  Product convergence of the independent reference
packages then gives the limiting law
\begin{equation}\label{tree:eq:full_contour_kernel_factor}
 W(x,a)\,\nu(dx,da)\,\mathcal B(x,a;dz)
       \prod_r K_r(x_r,dy_r).
\end{equation}
The factor $W$ need not factor across packages.  Both it and the
bridge law $\mathcal B$ depend only on the contour inputs, not on
the additional passage variables.  The complete contour and its
cut marks recover $(x,a,z)$, including the unused completion offsets.
Disintegration of Equation~\eqref{tree:eq:full_contour_kernel_factor}
therefore leaves exactly $\prod_rK_r(x_r,dy_r)$.

We remove the compact restrictions explicitly.  Exhaust the finite
cut-parameter domain by compact sets on which gap durations are
positive, then bound endpoints and offsets, and finally bound the
selected block paths in compact sets of the continuous-path topology.
Under the actual law these exclusions have probability tending to
zero, uniformly along the sequence, by the finite-block invariance
principle, tightness, and positivity of the finitely many limiting
cut gaps.  If $Q_n$ is the reference law and $P_n$ the actual law,
\[
 \E_{Q_n}[W_n\mathbf1_{A^c}]=P_n(A^c).
\]
Thus the same exhaustion controls the discarded density mass.
On each compact set the densities were uniformly bounded.  This
justifies the limiting density calculation and its normalization;
it does not assume a global bound for $W_n$.  A monotone-class
argument now gives the conditional identity for bounded measurable
tests.  Increasing bounded contour windows gives conditioning on
the entire bilateral contour.

\emph{Labels, reference densities, and overlapping tests.}
Before any limiting operation, the canonical graph map
restarts each block's clock at zero and its increments at zero.  Its
only noninternal numerical inputs are the specified forest differences.
It never uses the original integer cut index, an absolute starting
height, the starting time of a block, or any omitted jump count.
The global chronological permutation appears only in
Equation~\eqref{tree:eq:forest_density}.  Since this weight disappears
from the conditional output law in
Equation~\eqref{tree:eq:full_contour_kernel_factor}, that permutation is
not an additional parameter of $K_r$.

Changing a smooth positive reference density multiplies the joint
reference law by a function of the input variables alone. On a
compact parameter set the ratio of the lattice cell masses converges
uniformly to the continuous density ratio. The preceding test-function
argument, without bridge variables, therefore leaves the limiting
kernel $K_r$ unchanged. Exhaustion is justified under the replacement
law by tightness of the lattice approximations. The same argument
allows any smooth positive joint density for the finite forest-offset
vector within a package, independent of its strings: integrate over
the product lattice cells and use their common volume. Distinct
packages remain independent.

A permutation of block names is an exact relabeling of the finite
graph map. A change of forest basis is a unimodular integer linear
map $U$ of the offset vector, and recovering the same height
differences gives exactly the same graph. To compare the discretized
input laws, write $C=[-1/2,1/2)^d$ for the unit cell in this vector
space. The pushed-forward masses integrate the transformed smooth
joint density over cells $b+\epsilon_n UC$; its direct discretization
uses $b+\epsilon_n C$. Both cells have volume $\epsilon_n^d$ and
diameter tending to zero. Their masses divided by this volume
converge uniformly on compact sets to the same positive density.
Thus the ratios of the two masses tend uniformly to one there.
Both lattice laws are tight, as each is a vanishing displacement
of the same continuous input law, so their total variation distance
tends to zero. Applying the graph map preserves this bound even
when its output is discontinuous. Density invariance now identifies
the new-basis kernel with the transformed old kernel. This also
allows different positive reference densities for different finite
tests. Consequently a locally measurable enumeration of blocks
leaves no global label in the kernel.

For overlapping tests retain their joint output in one package and
apply the preceding argument to the union of their finite graphs.
The independently sampled blocks in this common experiment are always
the maximal consecutive runs of the union.  Adjacent pieces inside
one such run are recorded subwindows of a single block, with their
relative starts determined by that block's increments.  They are
never treated as separate independent blocks with a zero-duration
gap or a smooth density on their forced offset.  Thus all gaps to
which Equation~\eqref{tree:eq:forest_density} is applied remain positive.
If two forest choices are required, express their union of differences
in a single spanning forest, as above.  For disjoint spatial packages,
the block sets and their forests are disjoint, so their union is a
forest and the same completion calculation applies, even if their
chronological runs interleave.

\emph{Merging master blocks.}
If previously omitted
gaps are included, old separate strings can become windows in a new
master string.  Their offsets are then functions of that master
string; it would be invalid to restrict an arbitrary version of an
old kernel to this forced-offset relation.  Instead, start with the
old independent strings, smooth forest offsets, and their old test
outputs.  Insert the positive-duration old gaps as conditional bridge
strings and multiply by the exact density
Equation~\eqref{tree:eq:forest_density}.  The resulting concatenated string
has exactly the law of the new master string.  The old output is
unchanged, since its active edge-witness windows and band list are
unchanged.  Calculate any new union graph as an additional output
component on this same master string.  Applying the already-proved
density and bridge convergence to the old output gives precisely
Equation~\eqref{tree:eq:full_contour_kernel_factor}; hence its conditional
law given the entire new master contour is its original $K_r$.
This identity is proved before conditioning the master path, when
the gap endpoint differences have their ordinary densities.  There
is no conditioning of a preexisting almost-everywhere formula on a
singular slice.  If an old gap has duration zero, it was already
within a single old master block and the active-window identity is
exact without a density change.  For several old tests retain their
whole vector of outputs, together with the new test, before the
same operation.  Thus joint outputs, rather than filtered paths of
a larger graph, prove common-refinement compatibility.
We do not assert that a new union output is independent of the newly
observed bridge strings.  The calculation of the old conditional
marginal first integrates out that new output; the full output
vector is retained separately in the master experiment.

Extracting the countably many finite packages jointly preserves
these identities by the same compact density and bridge argument.  After
conditioning on the full contour and cuts they hold simultaneously
outside one null set, and hence also for a countably valued
contour-measurable selection.  This proves the lemma.
\end{proof}

\begin{proposition}[Spatial product kernels]\label{tree:prop:locality}
For every open coordinate domain $U$, the collection of all certificates
with compact range in $U$ has a regular conditional law
$\mathsf K_U(h|_U,\widehat h|_U;\cdot)$ depending only on the restrictions
of the two fields to $U$.  If $U_1,\ldots,U_m$ are disjoint open sets,
these collections are conditionally independent given the complete
fields, with respective conditional laws $\mathsf K_{U_i}$.
These laws are consistent under restriction and exhaustion.  They
also exist on the free-field patches obtained by local absolute
continuity in Proposition~\ref{tree:prop:background_locality}.
Here the local collection retains projected endpoints and ranges,
costs, and ordered-subpath incidences; the original global contour-time
labels have been forgotten.  Relative parametrizations inside locally
retained cells may be used during the construction.

In particular one can preserve a sampled collection on each of finitely
many disjoint patches and complete it to a sample of the whole
certificate law, provided that the preserved collections have the
specified product conditional law for the new fields on those patches.
This statement concerns conditional sampling; it makes no assertion
that the competitor satisfies a Weyl rule.
\end{proposition}

\begin{proof}
\emph{Finite charts with marks retained.}
Fix finitely many buffered charts in
each of the disjoint sets.  Refine and merge their instructions within
each set.  Different sets use disjoint cells.  Maximal runs from
different sets cannot share a cut: the shared point would belong to
both compactly contained cells.  Thus their distinct runs have
positive intervening time gaps.  Lemma~\ref{tree:lem:finite_kernel_transfer}
applies to these packages.

The reference kernel in that lemma has four kinds of input: relative
area durations, increment functions, internal relative cuts, and
forest height differences.  Lemma~\ref{tree:lem:chart_reconstruction}
reads all four from the oriented cells and boundary lengths in the
given patch.  The same local data give the embedding map that sends
the time-labelled ranges in the reference output to spatial ranges.
The conditional output law, after that map, is therefore a function
of the fields and marks \emph{in that patch}.  Lemma~\ref{tree:lem:finite_kernel_transfer}
also proved invariance under global block labels, reference densities,
and forest bases.  Thus replacing global cut indices by a local
enumeration of the cells does not discard a parameter on which the
kernel might still depend.  This explicit invariance is necessary;
conditional factorization given extra global labels would not by
itself prove the present assertion.

\emph{Conditioning on the fields.}
Realize the continuum fields by sampling their regular
conditional law given the peanosphere contour, independently of the
additional raw, time-labelled reference outputs and all time cuts
conditionally on that
contour.  This
is simply the prescribed mating-of-trees coupling, with any residual
embedding or imaginary-field randomness supplied at this stage.
Consequently the product identity from
Equation~\eqref{tree:eq:full_contour_kernel_factor} remains true on
conditioning additionally on these fields.  Its right-hand side
has just been shown measurable in the local fields and local marks.
Taking conditional expectation with respect to the complete fields
and marks gives the same product identity.  No determination of the
additional passage data by the continuum fields was used.

\emph{Exhaustion.}
Enumerate the countable instruction templates and
their rational positive-margin tests.  Within a patch, enumerate the
locally retained cells and all finite selections from them, using any
fixed measurable ordering of their embedded marked points.  Every
finite set of these tests has the conditional law just proved, with
the consistency supplied by the common-refinement part of
Lemma~\ref{tree:lem:finite_kernel_transfer}.  Hence they define a conditional
law on the countable product of compact certificate spaces.  Cylinder
events form a determining algebra.  The capture and realization
parts of Lemma~\ref{tree:lem:chart_reconstruction}, followed by an
exhaustion $K_j\Subset U$, identify the measurable collection generated
by this algebra with the local projected part of the original retained
array $\mathcal H$.  More explicitly, $\mathcal H_n$ and all the
charts were calculated from the same word and were extracted jointly.
Capture puts every compactly confined element of its limit in the
chart exhaustion; realization gives the reverse inclusion.  Both
statements include every finite number of ordered marks.  Thus no
new coupling between a metric and Poisson marks is chosen at the
limit.  The time cuts remain independent of the original array
and background, since their finite-resolution joint law was a
product and they were retained jointly.
They also identify the restriction maps for nested domains.  The
finite product identities on cylinder events extend by a monotone
class to the full local collections and then to countably many
subdomains when needed.

\emph{Removing the marks.}
Conditional on $(h,\widehat h,\eta,\mathcal H)$, the time cuts are still the
independent Poisson processes used at the start.  Since
$\eta_*(dt)=\mu_h$, the Poisson mapping theorem shows that the spatial
marks have their prescribed intensities proportional to $\mu_h$.
Their law consequently has that same form conditional on both fields,
and their restrictions to disjoint sets are independent under this
conditioning.  Nested refinements are obtained by independent
local superposition.  Integrating the marks out of the preceding
conditional product law therefore gives
\[
 \bigotimes_i\left(\int
  \mathsf K_{U_i}^{\mathrm{marks}}(h|_{U_i},\widehat h|_{U_i},\pi_i;\cdot)
     \,\operatorname{Pois}_{h|_{U_i}}(d\pi_i)\right).
\]
The factor in parentheses defines $\mathsf K_{U_i}$.  This proves
fields-only dependence and factorization, rather than merely a
mark-conditioned assertion.

\emph{Field changes and completion.}
On a relatively compact patch, change the field marginal by a locally
absolutely continuous density and retain the same conditional kernel.
The identities are preserved because the density is a function of
the conditioning variables.  Two such constructions agree on their
overlap by mutual absolute continuity and the established restriction
consistency.  An exhaustion yields the stated free-field versions.
All spaces used are standard Borel, so regular conditional laws exist.
Disintegrate the whole law over its restrictions to the preserved
patches to obtain the completion assertion.  If a recipe is subsequently
selected using the full sampled configuration, restriction to the
selection event gives a submeasure of this unrestricted construction;
it does not require a new conditional-independence claim.
\end{proof}

\subsection{Open margins, nondegeneracy, and changes of units}

The local kernels are now constructed. We next record the properties
needed when the laws themselves vary: inequalities pass through open
margins, the augmented construction survives further limits, and a
positive internal norm prevents all local costs from vanishing.

\begin{lemma}[Buffered tests]\label{tree:lem:certificate_buffered_tests}
Consider a convergent sequence in the certificate topology with a
uniform interval modulus.  Fix $K\Subset U$, a finite cost cutoff,
and compact endpoint sets in $K$.  The following implications hold.
\begin{enumerate}
\item An upper connection bound in the limit with a strict cost
margin and with certificates compactly in $U$ implies the same bound,
with a smaller margin, for the approximants.  Uniform bounds over
compact endpoint sets follow from finitely many witnesses and the
joining modulus.
\item A sequence of paths in $K$ violating a fixed lower traversal
bound, with bounded costs and endpoints in prescribed closed sets,
has a limiting certificate violating the corresponding weak bound
in $U$.  Any finite list of ordered crossings or subpaths can be
retained with it.
\item Suppose reference metrics $D_n$ are retained in the joint
topology and converge locally uniformly to a continuous length
metric $D$ inducing the ordinary topology.  Local inequalities
$F_n\leq C_nD_n$ along reference geodesics in arbitrary open margins,
and lower bounds on all certificates by $c_nD_n$, pass to the limit
when $C_n\to C<\infty$ and $c_n\to c\geq0$.
\end{enumerate}
These implications also hold under any fixed absolutely continuous
change of the retained continuum marginal.
\end{lemma}

\begin{proof}
For the first assertion, start with one limiting certificate.  Its
range has a positive distance from $U^c$, and its cost is strictly
below the bound, so its approximants retain both margins.  For a
compact endpoint family, continuity and uniform joining provide
neighborhoods of the endpoints in which one may use this certificate
with a prescribed small additional cost.  Finitely many such
neighborhoods suffice.  This gives the uniform finite-net assertion,
with no claim at $\partial U$.

For the second assertion, mark the specified crossing endpoints
before extraction.  Proposition~\ref{tree:prop:certificate_limit} gives
a converging marked certificate subsequence.  Closed endpoint
constraints and non-strict cost upper bounds persist in this
subsequence; its range is in $K\Subset U$.  This proves the required
contrapositive of the lower-bound transfer.

For the upper part of the third assertion, take a $D$-geodesic with
range compactly in an open margin.  Partition it into sufficiently
short segments.  The partition is finite.  For each segment,
continuity of $D$ and positive $D$-distance from its small surrounding
neighborhood to that neighborhood's outer boundary imply that a
$D_n$-geodesic between its endpoints stays in this neighborhood for
all sufficiently large $n$.  To justify existence and containment
one may work in a larger compact length-space chart: a competitor
of sufficiently small $D_n$-length beats every path reaching the
outer boundary, whose endpoint distance is bounded below by local
uniform convergence.  The sum of these $D_n$-distances converges to
the sum of the corresponding $D$-distances, which is the length of
the original geodesic.  Apply the assumed upper comparison on each
piece, concatenate, and extract.  The lower comparison follows
directly by applying it to approximants of a limiting certificate.
For a fixed change of marginal density $w$, approximate $w$ in
$L^1$ by bounded continuous functions of the retained continuum
data.  The convergence assertions pass for each such approximation,
and the $L^1$ error bounds the discarded probabilities uniformly.
\end{proof}

\begin{lemma}[Closure under further limits]\label{tree:lem:certificate_class_closed}
Augment a continuum certificate law by its original raw, time-labelled
array $\mathcal H$ and its projected local arrays as separate coordinates,
the common fixed continuum background, the independent
time cuts, and the laws of all finite independent reference packages
in Lemma~\ref{tree:lem:finite_kernel_transfer}.  The packages retain all
active-window tests, old and new outputs under master-block merging,
and the countably many fixed unit transforms.  Their continuum input
marginals are the fixed reference marginals $\nu_r$ in that lemma.
Suppose a sequence of such augmented laws has a common interval
modulus and the internal moment uniform integrability above.  Every
joint weak subsequential limit of these augmented data has the same
certificate, reference-transfer, reconstruction, and spatial
product-kernel properties.  In particular this class is closed
under the varying-law and varying-zoom extractions used below.
\end{lemma}

\begin{proof}
The input marginals of the reference experiments are fixed. This lets
us pass their density identities to a new limit, then recover the local
spatial collections from the retained raw arrays.

Write the $r$th reference law at stage $j$ as
$\nu_r(dx_r)K_{j,r}(x_r,dy_r)$.  The output spaces are compact,
the variables $y_r$ are the raw time-labelled outputs, and the input
marginals are fixed.  After joint extraction these
laws converge to $\nu_r(dx_r)K_r(x_r,dy_r)$.  For independent
packages, their product laws converge to the product of these
limits.  This is a statement about independent full experiments,
not a claim that arbitrary conditional independence is weakly closed.

For a fixed finite family, first restrict to compact input parameters
with positive gap durations.  The density $W$ and Brownian bridge
law in Equation~\eqref{tree:eq:full_contour_kernel_factor} are now the
same for every $j$.  On this compact set $W$ is bounded and
continuous; the bridge law is weakly continuous in the endpoints
and positive durations, by the bridge argument in
Lemma~\ref{tree:lem:finite_kernel_transfer}.  Integrating bounded
continuous bridge tests therefore preserves joint weak convergence
of the reference packages.  Hence that equation holds with the
new $K_r$.  The error in removing the compact restrictions is the
probability of the excluded parameters under the actual contour
law.  This law is fixed and does not depend on $j$ or on the
additional passage kernels.  Its exhaustion error tends uniformly
to zero.  Thus the full-contour factorization, label removal, and
forest-basis invariance all pass to the limit.

For merging, retain the entire vector of old and new test outputs
in the master experiment.  The old component is obtained at finite
resolution using the same active edge-witness windows.  The
positive-gap density-and-bridge construction near the end of
Lemma~\ref{tree:lem:finite_kernel_transfer} proves its conditional
marginal identity under observation of the new master string.
That construction uses the same continuous $W$ and bridge law,
so its identity passes by the preceding argument.  The output
coordinates and their deletion-of-marks maps are continuous maps
on the retained compact arrays.  Joint consistency of every
finite group of tests is consequently preserved.  No conclusion
is obtained by restricting a merged distance or by evaluating an
old kernel on a forced-offset diagonal.

First pass the fixed conditional coupling of the additional continuum
background with the contour, using only raw output tests.  For a
bounded background test, approximate its conditional expectation
given the contour by bounded continuous functions in $L^1$ of the
fixed contour marginal.  Multiplication by a bounded test of the
raw $\mathcal H$ or raw reference outputs keeps the error uniform
over $j$.  Joint weak convergence then preserves this conditional
sampling identity.  This argument is applied before projection;
projected outputs also depend on the background and are not claimed
conditionally independent of it.  It explains why the common full
background and its unit transforms must be retained, as in
Section~\ref{tree:sec:coherent-background}, rather than sampled separately at
successive scales.

Now project raw endpoints and ranges through the retained curve.
On each bounded time window, $(\eta,K)\mapsto\eta(K)$ is continuous
for uniform convergence of the curve and Hausdorff convergence of
the compact range; the same holds for every finite list of marks.
Thus the raw arrays and background determine their joint projected
limits.  The projected local part of the original $\mathcal H$
agrees with the chart exhaustion, with open margins:
on each compact $K\Subset U$, the proof of
Lemma~\ref{tree:lem:chart_reconstruction} selects finitely many bands
and exact minimum tests. Positive omitted-gap and spatial margins
persist in the contour/curve coupling after a slight enlargement
of confinement. Selected-minimum caps, which need not be strict,
remain exact comparisons inside the same retained finite test
outputs. Thus their treatment does not require continuity of the
graph map at an equality. Upper inclusion
uses approximating certificates, finite endpoint nets, and the
common joining modulus; lower inclusion uses bounded-cost marked
subpath extraction, as in
Lemma~\ref{tree:lem:certificate_buffered_tests}.  A countable exhaustion
in the compact, cost cutoff, number of marks, and positive margins
therefore preserves both inclusions.  This proves reconstruction
for the projected part of the original limiting array, rather than for a newly coupled
output.

Independence of the raw original array, together with its background,
from the time cuts passes as a product-law identity.
The locally measurable projection and the Poisson mapping theorem
now give Proposition~\ref{tree:prop:locality} by its same disintegration
argument.  Finally all fixed unit transforms are exact transforms
of the retained canonical experiments; their compatibility passes
by the same density calculation.  This proves every assertion.
\end{proof}

\begin{lemma}[A nonzero internal norm gives a nonzero local cost]
\label{tree:lem:nonzero_local_cost}
Let a sequence of discrete or already-continuum certificate laws
have the uniform interval modulus and uniform integrability of
internal endpoint $k$th powers from
Lemmas~\ref{tree:lem:moment_compactness} and~\ref{tree:lem:moment_ui}.
Suppose its internal endpoint costs on $[0,1]$ have a limiting
$k$th moment strictly greater than zero.  Then its limiting local
certificate costs cannot all vanish almost surely.

In particular, suppose reference metrics have a fixed continuum
law and local upper comparison constants $C_j\downarrow0$ along
such a sequence.  A joint subsequential limit cannot have a
positive internal unit-interval endpoint $k$th moment.  The assertion
allows different laws and different deterministic quantum-scale
zooms at each step of the sequence.
\end{lemma}

\begin{proof}
Let $Q=\eta([0,1])$.  Almost surely this is a disk with connected
interior and accessible endpoint tips, and the only preimages of
its interior are in $(0,1)$; see
Theorem~\ref{tree:thm:background_peanosphere} and
Equation~\eqref{tree:eq:cell_interior_preimages}.  There are times
$a_j\downarrow0$ and $b_j\uparrow1$ for which $\eta(a_j)$ and
$\eta(b_j)$ are interior points.  Indeed every initial and terminal
positive time interval contains points of the interior: its image
has positive area, while the boundary of the cell has zero area.
The times can be chosen from a countable dense set, since the
preimage of the interior is open.

Assume every compactly local cost vanishes.  Join $\eta(a_j)$ to
$\eta(b_j)$ by an ordinary continuous path in the cell interior.
Its compact range is covered by finitely many open sets compactly
in that interior.  Subdivide the path so each successive pair is
in one of these sets.  A zero local cost gives certificates of
arbitrarily small positive cost for these finitely many pairs.
Concatenation and the interior clause of
Lemma~\ref{tree:lem:chart_reconstruction} lift them to paths using only
the internal links of $[0,1]$.  Prepend and append the internal
short-time paths on $[0,a_j]$ and $[b_j,1]$.  The interval modulus
makes their costs tend to zero as $j\to\infty$.

To avoid any interchange of these choices with a limiting infimum,
first fix $j$ and a positive middle-path cost tolerance.  There
are finitely many certificates, all compactly in the cell interior.
They have actual approximating internal paths for all sufficiently
large resolutions.  The triangle inequality for the internal
endpoint distance therefore bounds its limit by the two endpoint
modulus terms plus this tolerance.  Send the tolerance to zero
and then $j$ to infinity.  The limiting internal endpoint cost is
zero almost surely.  Uniform integrability now makes its limiting
$k$th moment zero, contradicting the hypothesis.

For the last assertion, use local compactness of the reference
metric law and Lemma~\ref{tree:lem:certificate_buffered_tests}.  An upper
comparison constant tending to zero supplies zero-cost limiting
certificates along every sufficiently short reference geodesic.
A finite chain inside an open set gives zero local costs between
any two of its points; a countable dense set and the joining
modulus extend this to all pairs.  The first assertion is a
contradiction.  Its proof used only the common compactness, moment,
and interior-lifting properties, so it also applies to sequences
of continuum laws and changing zooms.
\end{proof}

It remains to identify how the construction changes under a coordinate
map or a change of quantum units. The latter may change the conditional
kernel and will be tracked explicitly.

Write $\phi^*$ for the transport of endpoints and ranges under a
conformal coordinate map $\phi:U\to V$; costs are unchanged.  For
$a>0$, write $S_a$ for multiplication of every cost, including
every marked-subpath cost, by $a$.  The imaginary-field pullback
is the usual imaginary-geometry transform
$\widehat h\circ\phi-\chi\arg\phi'$, taken modulo its angular
period.  Only its action on the directed curve and cells is used
here.

\begin{proposition}[Intrinsic transport and fixed quantum units]
\label{tree:prop:zoom_transport}
The local kernels of Proposition~\ref{tree:prop:locality} can be chosen
consistently under restriction so that the following law identities
hold.
\begin{enumerate}
\item For every fixed conformal coordinate map $\phi:U\to V$,
the pullback of the conditional certificate law on $V$ is the
conditional law on $U$ at the pulled-back fields
\[
 \widetilde h=h\circ\phi+Q\log|\phi'|,
 \qquad
 \widetilde{\widehat h}=\widehat h\circ\phi-\chi\arg\phi'.
\]
This is intrinsic coordinate covariance in law.
\item For fixed $v>0$, put $c_v=\gamma^{-1}\log v$ and define
the \emph{new} normalized law
\begin{equation}\label{tree:eq:level_kernel_definition}
 \mathsf K_U^{(v)}(g,\widehat g;\cdot)
  :=(S_{v^{-p}})_\#
     \mathsf K_U(g+c_v,\widehat g;\cdot).
\end{equation}
As a transformation of the joint law, this gives the real field its
translated marginal, namely the law of $h-c_v$, and preserves the
restriction and product-kernel properties on every $U$. On the buffered
reference patches of Proposition~\ref{tree:prop:background_locality},
local equivalence also defines it under the unchanged Gaussian
reference marginal. It need not equal $\mathsf K_U$.
For the established metric one has
\begin{equation}\label{tree:eq:reference_quantum_units}
             v^{-p}D_{g+c_v}=D_g.
\end{equation}
\item All fixed dyadic $v$ can be retained in one joint
extraction.  In the time-$2^j$, cost-$A_j$ experiment, the law
in Equation~\eqref{tree:eq:level_kernel_definition} uses time resolution
$2^{j-i}$ and cost denominator $A_j2^{-ip}$ when $v=2^{-i}$.
Consequently, if the level $j-i$ is also expressed in its own
$A_{j-i}$ units, its deterministic conversion factor is
\begin{equation}\label{tree:eq:zoom_normalization_factor}
 b_i=2^{ip}\lim_j\frac{A_{j-i}}{A_j},
\end{equation}
whenever this positive finite limit is part of the extraction.
\item For a fixed physical affine chart $\phi(z)=z_0+rz$, write
$h\circ\phi+Q\log r=g+c_v$.  The conditional law of the physical
certificate costs pulled back to this chart and divided by $v^p$
is exactly $\mathsf K_U^{(v)}(g,\widetilde{\widehat h};\cdot)$.
This identity remains valid on conditioning on exterior coarse
variables and then choosing a dyadic index measurably from those
variables, whenever the conditional interior field marginals are
absolutely continuous with respect to the local reference marginals.
\end{enumerate}
These are identities at each fixed index, before any limit as $v\to0$.
Changes of density under a fixed reference marginal always use patches
whose full dependency region avoids its pinning circle. No uniform
absolute-continuity bound for unbounded constant shifts is asserted.
\end{proposition}

\begin{proof}
An intrinsic chart records quantum areas, boundary length increments,
directed arc roles, and shared seams.  These quantities are unchanged
by the quantum and imaginary-geometry coordinate transformations.
The instruction graph consequently stays exactly the same, while
the map taking its time-labelled ranges to the coordinate plane is
composed with $\phi^{-1}$.  Reference-kernel independence of labels
and parameter densities, followed by local reconstruction, therefore
gives the first assertion on every finite chart.  Exhaustion gives
the full local collection.  This is a coordinate identity for the
constructed intrinsic experiment, rather than an application of
metric uniqueness.  Coordinate covariance of $D_h$ is separately
the known theorem \cite[Theorem~1.3]{GwynneMillerConformal}.

At the level of joint laws, transform the real field by $h\mapsto h-c_v$
and every cost by $S_{v^{-p}}$. Disintegration gives
Equation~\eqref{tree:eq:level_kernel_definition} under the translated
field marginal, on any domain $U$. Conditional factorization and
restriction consistency pass through this deterministic transformation.

To compare with an unchanged Gaussian reference marginal, fix the
pinning circle outside the entire bounded region used by the tests,
their dependency buffers, and the shift supports. Extend the constant
$c_v$ on a test patch to a smooth function supported away from that
circle. The extension has zero average on the pinning circle and is
a Cameron--Martin direction of this same normalized field law. Hence
the shifted restriction is locally equivalent to the reference
restriction, so the right side of
Equation~\eqref{tree:eq:level_kernel_definition} is defined almost surely
there. This argument is used only on those patches; a domain containing
the entire pinning circle would observe its fixed average. The
established Weyl rule gives
$D_{g+c_v}=e^{\xi c_v}D_g=v^pD_g$, proving
Equation~\eqref{tree:eq:reference_quantum_units}.

We prove that this definition is the required time zoom, since a
definition of unrelated kernels would not suffice.  Adding $c_v$
multiplies the quantum area of every fixed cell by $v$ and its
quantum boundary lengths by $\sqrt v$.  Accordingly a canonical
block with duration $\ell$, increment path $z(t)$, and offsets $o$
is transformed into
\begin{equation}\label{tree:eq:block_unit_transform}
 \ell\mapsto v\ell,\qquad
 z(t)\mapsto\sqrt v\,z(t/v),\qquad
 o\mapsto\sqrt v\,o.
\end{equation}
Internal cuts transform with the duration. Anchored band endpoints
$\alpha,\beta$ are multiplied by $\sqrt v$, as are their anchor
heights and selected minima. These are the transformed instruction
templates retained at the start of this section. Lattice-cell
discretization of offsets commutes exactly with this scalar change
of variables. At finite resolution,
a rate-$n$ Poissonized string in a block of duration $v\ell$ has
exactly the same distribution, after the time change, as a
rate-$nv$ string in a block of duration $\ell$.  Its height lattice
also agrees: the latter lattice, multiplied by $\sqrt v$, is
$\sqrt v(nv)^{-1/2}\Z=n^{-1/2}\Z$.  The matching tests, full height
strips, and selected-minimum comparisons are the same integer
tests on the same string.  Thus this is an exact finite-experiment
identity, including every path and every marked subpath.  Changing
the positive duration and offset sampling densities under
Equation~\eqref{tree:eq:block_unit_transform} does not alter the
conditional kernel, by Lemma~\ref{tree:lem:finite_kernel_transfer}.
Taking its joint limit proves the asserted identity with time
resolution $nv$ and denominator $B_nv^p$.  Substituting
$n=2^j$, $B_n=A_j$, and $v=2^{-i}$ gives
Equation~\eqref{tree:eq:zoom_normalization_factor}.

The continuum background is transported at the same time.  If
$(h,\widehat h,\eta,Z)$ is the original realization, normalizing
the physical area-$v$ patch on the same abstract surface means
\begin{equation}\label{tree:eq:simultaneous_background_zoom}
 h^{[v]}=h-c_v,\qquad
 \eta^{[v]}(t)=\eta(vt),\qquad
 Z^{[v]}(t)=v^{-1/2}Z(vt),\qquad
 F^{[v]}=v^{-p}F.
\end{equation}
In particular $D_{h^{[v]}}=v^{-p}D_h$.  If a canonical cone
embedding is desired, use its measurable conformal re-embedding
and the first assertion.  The Brownian scaling and cone embedding
laws are the standard ones in Theorem~\ref{tree:thm:background_peanosphere}.
All fixed $v$ use these transformations of a single background and
the same microscopic experiment.  Recording countably many fixed
$v$ before any further subsequence retains their compatibility.
Equivalently, one may use a common Brownian/walk coupling with
uniformly negligible rescaled contour error and apply exactly
the transformations in Equation~\eqref{tree:eq:simultaneous_background_zoom}.

For the final assertion, the first part says that the pulled-back
physical law is $\mathsf K_U(g+c_v,\widetilde{\widehat h};\cdot)$.
Multiplying its costs by $v^{-p}$ is precisely
Equation~\eqref{tree:eq:level_kernel_definition}.  One common null set
suffices for the countable fixed dyadic indices and fixed coordinate
charts used in the proof.  Given exterior coarse data, apply these
identities on the conditional interior field space.  Absolute
continuity of its marginal makes the same versions valid there;
selection of a dyadic index is a countable measurable partition.
For random harmonic residual data, the argument applies to the
whole conditional field distribution.  It does not substitute
uncountably many deterministic field values into arbitrarily
chosen versions.  This proves the last assertion and the stated
order of operations.
\end{proof}

\subsection{Coherent backgrounds and the law topology}
\label{tree:sec:coherent-background}

We next put all scales in one space of observations.  Its purpose is
to retain the same continuum background under time changes and to
make the later success tests stable under extraction.

\paragraph{Retained observations.}
Enumerate the
bounded rational time windows, the finite Poisson/rational chart
instructions of Section~\ref{tree:sec:certificates}, their rational
open margins and cost cutoffs, and all finite lists of subpath
marks.  For each instruction use its compact closed-certificate
hyperspace from Proposition~\ref{tree:prop:certificate_limit}.  Record
also the internal endpoint variables in $[0,\infty]$, the
continuous contour on compact windows, and the continuum
background used in Proposition~\ref{tree:prop:zoom_transport}.
The background includes the curve and reference metric in local
uniform topology and the fields in their distribution topologies.  It is
given its usual standard Polish realization.  For every bounded
connected rational coordinate domain $W$ and every compact
rational exhaustion set $K\Subset W$, retain also
$D_h(\cdot,\cdot;W)|_{K\times K}$ in uniform topology.
These internal-reference-metric observations have a fixed
joint marginal; retaining them does not assert continuity
of the map from a field or an ambient metric to its internal
metric.  We also retain,
as deterministic coordinates, the joint laws of every auxiliary
finite reference experiment in
Lemma~\ref{tree:lem:finite_kernel_transfer}: the canonical input
and complete marked output law for each package, all finite
lists of independent packages, and their common refinements.
Use the same prescribed smooth input densities at every scale.
Include the countably many fixed quantum-unit transforms of
these templates.  At finite resolution the input marginals
depend on the random-walk resolution and offset lattice;
the invariance principle and the lattice-cell approximation of
forest offsets make them converge to the prescribed continuum input
marginals.  These limiting
input marginals are the same for every continuum law in a
subsequent varying-law extraction.  Each output space is a
compact hyperspace, so input tightness gives precompactness
of the finite-resolution weak-law coordinates as well.
They are auxiliary reference
laws, not extra samples from the actual contour.  Their
compatibility with the actual atlas, including the
density-and-bridge identities, is retained by
Lemma~\ref{tree:lem:certificate_class_closed}.

A bounded product
metric, with coordinate weights $2^{-r}$, makes this a separable
metric data space $\mathcal E$; use the closed compatible-data
subspace when a refinement identity is imposed.  We equip its
laws with a metric $\mathbf d$ for weak convergence, for example
the bounded Lipschitz metric for this bounded product metric.
Equivalently, this topology is specified by integrals against a
fixed countable convergence-determining family of bounded
continuous functions.  Thus all tests below refer to one
specified countable law topology.

The raw certificate coordinates are the time-labelled compact
spaces in Equation~\eqref{tree:eq:marked_certificate_space}.  Discrete
grid times are points of the same compact time window, so a
finite-resolution law belongs to this ambient space without
satisfying the Brownian quotient identities.  All choices of
representatives are retained.  Spatial projection is applied
only after adjoining the continuum background; the vanishing
projected diameter of individual links supplies its limiting
compatibility.

\paragraph{One Brownian coupling for all scales.}
Take independent
two-sided Brownian motions $W_1,W_2$ and an independent sequence
of fair coordinate choices.  For each coordinate, read successive
simple-walk increments from the signs of Brownian exits from the
interval of radius one about its current position.  Use the
coordinate choices to interleave the two resulting walks; apply
the same construction independently on the negative half-line.
This gives exactly the bilateral axial walk.  The successive
Brownian exit durations are independent, have mean one, and
satisfy the law of large numbers.  The numbers of each coordinate
choice likewise have asymptotic frequency $1/2$.  Hence the
Brownian clock used by coordinate $a$ at walk time $Nt$, divided
by $N$, converges uniformly on bounded $t$-intervals to $t/2$,
in probability.  The Brownian modulus of continuity after
diffusive rescaling implies
\begin{equation}
 \begin{gathered}
 \sup_{|t|\leq T}N^{-1/2}
 \bigl|Z^{\rm disc}(\lfloor Nt\rfloor)-B(Nt)\bigr|
 \longrightarrow0\quad\text{in probability},\\
 B(t)=(W_1(t/2),W_2(t/2)),
 \end{gathered}
 \label{tree:eq:bootstrap_common_coupling}
\end{equation}
for every fixed $T<\infty$.  The interpolation error is at most
one step.  This argument only uses the strong Markov property,
the exit-time mean, and the law of large numbers.

\paragraph{Attach and transport the continuum background.}
Let $\nu$ be the law of the full bilateral Brownian contour.
Write $\nu(db)\mathsf Q(b,dc)$ for its standard joint law with
the full continuum background $\mathcal C$, including any
residual imaginary-field and embedding data.  This background
excludes the microscopic coordinate schedule, time cuts, and
additional passage data.  If
$S_vb(t)=v^{-1/2}b(vt)$, its intrinsic quantum-unit transform is
\[
 T_v\mathcal C=(h-c_v,\widehat h,\eta(v\,\cdot),v^{-p}D_h),
 \qquad c_v=\gamma^{-1}\log v,
\]
followed by the prescribed common reembedding of its components.
The decorated-cone law is invariant under $(S_v,T_v)$ by
Equation~\eqref{tree:eq:background_joint_zoom}.
Here the imaginary field has the uniform phase modulo its
angular period specified in
Proposition~\ref{tree:prop:background_locality}.  Under a positive
circle-average dilation, its GFF modulo constants is invariant,
and the change of anchoring constant is absorbed by that
uniform phase.  The dilation is determined by the real field,
which is independent of the imaginary field before the curve
is drawn.  The curve then transforms by imaginary-geometry
covariance, and the reference metric by its established
covariance and constant-shift rule.  Thus the invariance is
for the full background, including the extra imaginary
randomness.  A measurable representative is chosen from the
intrinsic decorated surface and never from passage data.

Disintegration of this joint invariance gives
\begin{equation}
 \mathsf Q(S_vb,\cdot)=(T_v)_\#\mathsf Q(b,\cdot)
 \quad\text{for $\nu$-almost every $b$}.
 \label{tree:eq:bootstrap_background_kernel_equivariance}
\end{equation}
For the countable dyadic group, take one full-measure set on
which these identities and their compositions all hold.  Draw
$\mathcal C$ once from $\mathsf Q(B,\cdot)$, independently of
the coordinate schedule and all time-cut processes conditionally
on $B$.  Put $B_j=S_{2^j}B$ and $\mathcal C_j=T_{2^j}\mathcal C$.
The background at $j-i$ is exactly the time-$2^{-i}$ transport
of the background at $j$, intrinsically.  The discrete
instruction data obey the same identity, with the extra
deterministic cost conversion $A_{j-i}/A_j$.

Let $Y_j$ denote the raw, unprojected time-labelled passage
data and time cuts at scale $j$.  Since $S_{2^j}$ is invertible
on the full bilateral contour, $B_j$ determines $B$.
The conditional sampling and
\eqref{tree:eq:bootstrap_background_kernel_equivariance} give the
exact finite-resolution identity
\begin{equation}
 \Law(\mathcal C_j\mid B_j,Y_j)=\mathsf Q(B_j,\cdot).
 \label{tree:eq:bootstrap_background_conditional}
\end{equation}
One must use raw data here: their projected spatial ranges
also depend on $\mathcal C_j$ and are adjoined afterward.

\paragraph{Preserve the conditional law under extraction.}
We verify that \eqref{tree:eq:bootstrap_background_conditional}
passes through extraction despite possible discontinuity of
$\mathsf Q$.  Suppose $(B_j,Y_j)\Rightarrow(B_\infty,Y_\infty)$.
The $B_j$ marginal is always $\nu$, and the $\mathcal C_j$
marginal is always $\nu\mathsf Q$, so the joint triples are
tight.  For bounded continuous $f$ and $g$, write
$\mathsf Qg(b)=\int g(c)\mathsf Q(b,dc)$.  Then
\[
 \E[f(B_j,Y_j)g(\mathcal C_j)]
 =\E[f(B_j,Y_j)\mathsf Qg(B_j)].
\]
Approximate $\mathsf Qg$ in $L^1(\nu)$ by bounded continuous
functions with a common bound.  The replacement error is at
most $\|f\|_\infty$ times this $L^1$ error, uniformly in $j$
and also in the limit.  Weak convergence applies to each
continuous approximation.  Removing the approximation yields
\begin{equation}
 \Law(\mathcal C_\infty\mid B_\infty,Y_\infty)
       =\mathsf Q(B_\infty,\cdot).
 \label{tree:eq:bootstrap_background_closure}
\end{equation}
Products of continuous tests determine the joint law, giving
the stated conditional identity.  Thus conditioning the
limiting reference experiments additionally on the continuum
fields does not change their passage kernels.  This is a
fixed-kernel argument; closure when the passage kernels vary
uses the retained reference-law coordinates and
Lemma~\ref{tree:lem:certificate_class_closed} instead.

\paragraph{The law at scale $j$.}
Write $\mu_j$ for the law of these complete data in time units
$2^j$ and cost units $A_j$.  Equation~\eqref{tree:eq:bootstrap_common_coupling}
and Proposition~\ref{tree:prop:certificate_limit} give the required
joint extractions.  Recording the continuum background avoids
asserting that an arbitrary measurable functional of a contour
is continuous.  Its joint marginal is fixed, and its contour
equals the discrete contour limit.  The exact transport
identities just established identify all fixed-offset limits
with the corresponding zooms in a common extraction.

Fixed local changes of density can be included in this topology
without a new compactness assumption.  Truncate each density and
approximate it in $L^1$ of the fixed background marginal by bounded
continuous functions.  Weak convergence applies to the
approximations, and their $L^1$ errors are uniform over the
coupled laws.  Use a countable list of the buffered domains and
approximations required for the shell tests.  In particular, the
local free-field tests below are tests of the same kernels
identified in Proposition~\ref{tree:prop:locality}.


